% Pursuing Stacks (cotes 134-2 to 134-8): archive cote 134-3, site page 45 (PDF page 46), typescript page 45.
% typescript. Draft compared: ps2.tex:2225-2293.
% AI-transcribed by Claude Opus 5.5 (claude-opus-5-5, adaptive thinking, effort high), Message Batches API, batch ps1 (msgbatch_01Rv5hCaNTtUt4JHx1yqoHPL), 2026-11-12; not yet checked by eye.
\scanpage{134-3/45}
\mspage{45}
for homotopy types,  for  instance set-theoretic realisations of some speci-%
fied algebraic structure species, or topological spaces,  and the like,
do admit arbitrary direct and inverse limits,  and surely \xout{\unsure{finite}} sums and
products therefore,  so that the question of commutation of the localizati-%
on functor \inserted{to these,} arises indeed and is felt to be important.  Possibly so impor-%
tant even,  that the notion I introduced of a ``modelizer'' should take this
into account,  and be strengthened to the effect that the canonical functor
$M\longrightarrow$(Hot)  should commute at least to finite sums and products,  and possib-%
ly even to infinite ones (\xout{\unsure{this}} \inserted[typed]{whether} \inserted{the latter} will have to be looked up with care).  I'll
admit \inserted{provisionally,} that in (Hot),  finite sums and products can be described in terms of
the corresponding operations in (Cat),  namely that the canonical functor
(going with our very definition of (Hot) as a localization of (Cat))
\[ (\mathrm{Cat}) \longrightarrow (\mathrm{Hot}) \]
commutes with finite (presumably even infinite) sums and products.\hand{(\unsure{35})} This is
indeed reasonable,  in view of the fact that the Nerve functor
\[ (\mathrm{Cat}) \longrightarrow \hand{\Delta}^{\wedge} = (\mathrm{Ss\ sets}) \]
does commute to sums and products.

In the case of a test category\xout{,} A and the corresponding elementary mode-%
lizer $A^{\wedge}$,  the corresponding localizing functor is the compositum
\[ A^{\wedge} \xrightarrow{\ iA\ } (\mathrm{Cat}) \longrightarrow (\mathrm{Hot}) \ , \]
which therefore commutes with finite (presumably even infinite) sums automa-%
tically,  because $\hand{i}_A$ does.  Commutation with finite \uline{products} though does
not look automatic.  The \xout{fact that} property of commutation with final
elements is OK and is nothing but condition T 1),  which we saw is necessa-%
ry for A to be a test category.  Thus remains the question of commutation
with binary products,  which boils down to the following condition,  for
any two elements  F and G in $A^{\wedge}$:
\[ i_A(F\mathrm{x}G) \longrightarrow i_A(F)\mathrm{x}i_A(G) \quad \text{should be a weak equivalence} \]
i.e.
\[ (*) \qquad A_{/F\mathrm{x}G} \longrightarrow A_{/F}\ \mathrm{x}\ A_{/G} \quad \text{a weak equivalence .} \]
\xout{a weak equivalence} This now implies condition T 2),  as we see taking
$F = a$ , $G = b$ ,  \xout{and} in which case the condition \xout{\illegible{2}} \inserted[typed]{(*) just means} asphericity
of  $a\mathrm{x}b$ in $A^{\wedge}$ namely  T 2 .  To be happy,  we have still to show that
conversely,  (T 2) implies (*)  for any F,G.  As usual,  it implies even the
stronger condition that the functor in question is aspheric,  which by the
standar\over{t}{d} criterion  (page 38)  just means that the categories $A_{/a\mathrm{x}b}$
(for a,b in A, \inserted{and,} such \inserted{moreover,} that F(a) and G(b) non empty, \inserted{but} never mind) are aspheric.

Everything turns out just perfect  -  it seems worthwhile to summarize
it in one theorem:
